% (C) 2026 Ralf Stephan, in collaboration with Claude Code.
% Released under CC0 1.0 Universal (public-domain dedication).
% See https://creativecommons.org/publicdomain/zero/1.0/

\documentclass[11pt]{article}
\usepackage[T1]{fontenc}
\usepackage[utf8]{inputenc}
\usepackage{lmodern}
\usepackage[margin=2.5cm]{geometry}
\usepackage{amsmath,amssymb}
\usepackage{booktabs,tabularx}
\usepackage{braket}
\usepackage[hidelinks]{hyperref}

% Lean names and paths: verbatim, breakable at . _ /
\DeclareUrlCommand\lean{\urlstyle{tt}}
\def\UrlBreaks{\do\.\do\_\do\/}
\newcolumntype{L}{>{\raggedright\arraybackslash}X}
\newcolumntype{B}{>{\raggedright\arraybackslash}p{0.42\textwidth}}
\emergencystretch=6em

\title{Lean formalization of an Assumptions of Physics technical brief:\\ log}
\author{}
\date{}

\begin{document}
\maketitle

\noindent
This log records which statements of the technical brief in this directory were checked in Lean,
where the formal version differs from the brief, and any problems found in the brief.
The brief is available at \url{https://assumptionsofphysics.org/resources/briefs}.

\section*{Setup}

\begin{itemize}
\item Lean \texttt{v4.36.0-rc1}; Mathlib at commit \texttt{2cabef97}; physlib (\texttt{master})
  and TauCeti (commit \texttt{eb4dd24}, on the same Mathlib) as further dependencies.
\item Build: \texttt{lake build} in the project root.
\item Every file is checked with \texttt{\#print axioms} on its main results. A result is
  recorded here only if it depends on no axioms beyond \lean{propext},
  \lean{Classical.choice} and \lean{Quot.sound}, and so has no \lean{sorry}.
\item Imports are kept to the Mathlib modules actually needed.
\item Material that belongs in Mathlib is kept in \lean{ForMathlib/}, in Mathlib's directory
  layout (library \lean{ForMathlib}), and imported by the brief files.
\end{itemize}


%======================================================================

\section{Schr\"odinger equation is equivalent to information conservation}

\paragraph{Files.}\leavevmode
\begin{itemize}
\item File: \lean{briefs/UnitaryEntropy/UnitaryEntropy.lean}
\item Library: \lean{UnitaryEntropy}
\item Imports: \lean{Mathlib.Analysis.SpecialFunctions.BinaryEntropy}, the Wigner library files \lean{EqualMixtures.lean},
  \lean{Kadison.lean} and \lean{Differentiable.lean} (which import the rest of\newline
  \lean{ForPhyslib/QuantumMechanics/Wigner}), and two new \lean{ForMathlib} files:\newline
  \lean{ForMathlib/Analysis/Calculus/Deriv/Star.lean} (in a normed star algebra, the derivative
  of a path of unitaries through $1$ is skew-adjoint, and $U^\dagger U$ is constant when
  $U' = XU$ with $X$ skew-adjoint),\newline
  \lean{ForMathlib/Analysis/InnerProductSpace/Adjoint.lean} (in finite dimension,
  $U^\dagger U = 1$ already makes $U$ unitary)
\item Wigner library, in \lean{ForPhyslib/QuantumMechanics/Wigner/}: \lean{Density} (density
  operators, pure states \lean{proj}, equal mixtures \lean{mix}; moved there from this brief),
  \lean{Basic},
  \lean{Normalize}, \lean{Rigidity}, \lean{Theorem} (Wigner's theorem without bijectivity, after
  Geh\'er, now in any dimension, see Section~\ref{sec:unitary-entropy-any-dim}; uniqueness up to
  one phase in finite dimension), \lean{Projection} (the same for
  maps on rank-one projections), \lean{Continuity} (Bargmann's triple product excludes the
  anti-unitary case near $t_0$), \lean{Differentiable} (explicit, differentiable choice of the
  phase). Its \lean{ForMathlib} helpers are \lean{Analysis/Complex/ThreeCircles.lean},
  \lean{Analysis/InnerProductSpace/PhaseUnitary.lean} and
  \lean{Analysis/InnerProductSpace/LinearMap.lean} (conjugate-linear isometries conjugate the
  inner product; polarization of $\ket{x}\bra{y}$). Coordinate conjugation in an orthonormal
  basis came from TauCeti (\lean{TauCeti.Analysis.InnerProductSpace.Conjugation}) until the
  library moved to Hilbert bases.
\end{itemize}

\paragraph{Setting.} $E$ is a finite-dimensional complex Hilbert space, and operators are
continuous linear maps $E \to E$. A density operator (\lean{IsDensity}) is a positive operator of
trace $1$. The pure state of a unit vector $\psi$ is \lean{proj}~$\psi = \ket{\psi}\bra{\psi}$
(Mathlib's \lean{rankOne}), and \lean{mix}~$\rho\,\sigma = \tfrac12\rho + \tfrac12\sigma$. The
von Neumann entropy \lean{vonNeumann} is $\sum -\mu \log \mu$ over the roots $\mu$ of the
characteristic polynomial. For self-adjoint $\rho$ these are its eigenvalues
(\lean{vonNeumann_eq_sum}), and the definition makes invariance under conjugation immediate.
\lean{born}~$\psi\,\phi = |\braket{\psi|\phi}|^2/(\|\psi\|^2\|\phi\|^2)$, and
\lean{mixEntropy}~$p = H_2((1+\sqrt p)/2)$ with $H_2$ Mathlib's \lean{binEntropy}. The
Schr\"odinger equation for a propagator is $U'(t) = (i\hbar)^{-1} H(t)\,U(t)$ with $H(t)$
self-adjoint.

\paragraph{Results.}
\begin{center}
\begin{tabularx}{\textwidth}{@{}BL@{}}
\toprule
Brief & Lean \\
\midrule
Corollary 1, entropy of $\tfrac12\rho_\psi + \tfrac12\rho_\phi$ & \lean{vonNeumann_mix_proj},
  \lean{mixEntropy_eq} \\
Corollary 1, strictly decreasing & \lean{mixEntropy_strictAntiOn} (on $[0,1]$) \\
Corollary 1, $p(\phi|\psi) = \lambda$ iff $S = \ldots$ & \lean{born_eq_iff} (for
  $\lambda \in [0,1]$) \\
Theorem 2 ($\Rightarrow$), unitary propagator & \lean{star_mul_self_of_schrodinger},
  \lean{ContinuousLinearMap.mem_unitary_of_star_mul_self} \\
Theorem 2 ($\Rightarrow$), existence for constant $H$ & \lean{hasDerivAt_exp_schrodinger}:
  $U(t) = \exp(t(i\hbar)^{-1}H)$ \\
Theorem 2 ($\Rightarrow$), DR-INFO & \lean{vonNeumann_unitary_conj},
  \lean{IsDensity.unitary_conj}, \lean{schrodinger_conserves_entropy} \\
Theorem 2 ($\Leftarrow$), Step 1 & \lean{vonNeumann_eq_zero_iff}: $S(\rho) = 0$ iff $\rho$ is
  pure \\
Theorem 2 ($\Leftarrow$), Steps 1--2 up to Wigner & \lean{pure_and_born_of_preserves},
  \lean{proj_and_trace_of_preserves} \\
Theorem 2 ($\Leftarrow$), Step 2, Wigner's theorem & \lean{Wigner.wigner_rankOne} (from
  \lean{Wigner.wigner}) \\
Theorem 2 ($\Leftarrow$), anti-unitary case excluded & \lean{Wigner.eventually_exists_unitary},
  \lean{PhaseEqOnSphere.false_of_unitary_of_antiunitary} \\
Theorem 2 ($\Leftarrow$), differentiable $U(t)$ & \lean{Wigner.exists_unitary_differentiableAt} \\
Theorem 2 ($\Leftarrow$), Step 3 & \lean{exists_hamiltonian} (from
  \lean{star_eq_neg_of_hasDerivAt_of_mem_unitary}), for $\hbar \ne 0$ \\
Theorem 2 ($\Leftarrow$), mixed states & \lean{Wigner.affineOnDensities_of_map_mix},
  \lean{Wigner.AffineOnDensities.eq_conjOp}, \lean{Wigner.mul_mul_star_eq_conjOp} \\
Theorem 2 ($\Leftarrow$), complete & \lean{schrodinger_of_preserves} \\
\bottomrule
\end{tabularx}
\end{center}

\paragraph{Differences from the brief.}
\begin{itemize}
\item This part is finite-dimensional; the brief has a separable Hilbert space. The calculus
  results (\lean{star_mul_self_of_schrodinger}, \lean{exists_hamiltonian}) hold in any Hilbert
  space for bounded $H(t)$. Corollary~1 and Theorem~2 through Wigner's theorem are redone in
  any dimension in Section~\ref{sec:unitary-entropy-any-dim}.
\item Corollary~1 is proved without diagonalizing the mixture. $\rho = \tfrac12\rho_\psi +
  \tfrac12\rho_\phi$ satisfies $\rho^3 = \rho^2 - \frac{1-\lambda}{4}\rho$
  (\lean{mix_proj_cubic}), so every eigenvalue is $0$ or $(1 \pm \sqrt\lambda)/2$. On these three
  points $-\mu\log\mu = \alpha\mu + \beta\mu^2$ for suitable $\alpha, \beta$, and
  $\operatorname{tr}\rho = 1$, $\operatorname{tr}\rho^2 = (1+\lambda)/2$ fix the sum. This also
  covers the degenerate cases $\lambda = 0$ (double eigenvalue $\tfrac12$) and $\lambda = 1$
  ($\rho$ pure).
\item [B]EVMIX is used only for equal mixtures, which is all Step~2 needs. For the
  ($\Rightarrow$) direction, $\rho \mapsto U\rho U^\dagger$ is linear, so it preserves all mixtures.
\item DR-INFO is $S(\rho_I) = S(\rho_F) = S(\rho_P)$. The reconstructed past
  $U^\dagger\rho_F U$ is $\rho_I$ itself, which is also proved.
\item The ($\Rightarrow$) direction shows that every propagator solving the equation is unitary.
  Existence is shown for constant $H$ only; for time-dependent $H(t)$ it is the linear ODE
  existence theorem, which the brief also assumes.
\item \lean{schrodinger_of_preserves} assumes $\Phi_{t_0} = \mathrm{id}$ and differentiability of
  $t \mapsto \Phi_t(\rho)$ at $t_0$ only on density operators, one $\rho$ at a time. It concludes
  that $\Phi_t(\ket{\psi}\bra{\psi}) = \ket{U(t)\psi}\bra{U(t)\psi}$, and
  $\Phi_t(\rho) = U(t)\rho U(t)^\dagger$ for every density operator $\rho$, for $t$ near $t_0$, with
  $U(t)$ unitary for every $t$, $U(t_0) = 1$ and $U'(t_0) = (i\hbar)^{-1} H U(t_0)$ for a
  self-adjoint $H$. This is the Schr\"odinger equation at $t_0$, which is what EVDIFF at $t_0$
  gives. For a two-parameter evolution, one propagator solves it at every time
  (Section~\ref{sec:unitary-entropy-every-time}).
\item Wigner's theorem is used in its non-bijective form (Geh\'er), so the bijectivity of the
  correspondence of pure states via the reconstructed past (Step~1) is not needed.
\item $\Phi_t(\rho) = U(t)\rho U(t)^\dagger$ for mixed $\rho$ follows from the pure states once
  $\Phi_t$ is affine, by the spectral decomposition of $\rho$
  (\lean{Wigner.AffineOnDensities.eq_conjOp}). Affinity is not an extra assumption: a map of the
  density operators that preserves equal mixtures preserves all convex combinations of two of
  them (\lean{Wigner.affineOnDensities_of_map_mix}), since it is bounded and preserves midpoints.
\end{itemize}

\paragraph{Problems found in the brief.} No errors. Implicit points:
\begin{itemize}
\item The ``iff'' in Corollary~1 needs $\lambda \in [0,1]$, where the entropy is strictly
  decreasing in $\lambda$; for other $\lambda$, $\sqrt\lambda$ and the logarithms are not
  meaningful.
\item Step~3 assumes a norm-differentiable unitary family, which forces a bounded generator. On
  an infinite-dimensional Hilbert space it therefore covers bounded Hamiltonians only, and even
  then EVDIFF does not give such a family
  (Section~\ref{sec:unitary-entropy-norm}). The brief ``Multiple
  equivalent characterizations of the Schr\"odinger equation'' treats the general case with
  Stone's theorem.
\item In Step~3, $\hbar \ne 0$ is needed to write $A = H/(i\hbar)$.
\item Wigner's theorem gives each $U(t)$ only up to a phase, and Step~3 differentiates $U(t)$. A
  differentiable choice exists. Extend $\Phi_t$ to all $\ket{x}\bra{y}$ by polarization, so
  that it sends $\ket{x}\bra{y}$ to $\ket{Ux}\bra{Uy}$, and fix an orthonormal basis $(b_k)$ and
  an index $i$. Then
  \[ U(t) = \bigl(\operatorname{Re}\braket{b_i|\Phi_t(\ket{b_i}\bra{b_i})\,b_i}\bigr)^{-1/2}
    \sum_k \Phi_t(\ket{b_k}\bra{b_i})\,\ket{b_i}\bra{b_k} \]
  is the unitary with $\braket{b_i|U(t)b_i} > 0$. It is a finite expression in values of $\Phi_t$ on pure states, so
  it is differentiable at $t_0$ whenever these are.
\item ``Joint continuity and $\Phi_{t_0} = I$ exclude the anti-unitary case'' needs only continuity
  at $t_0$ of $t \mapsto \Phi_t(P)$ for three pure states $P$: the imaginary part of Bargmann's
  triple product $\operatorname{tr}(\Phi_t(P_1)\Phi_t(P_2)\Phi_t(P_3))$ is $\tfrac14$ at $t_0$ and
  would be $-\tfrac14$ for an anti-unitary. In dimension $1$ there is nothing to exclude.
\end{itemize}

\subsection{The equation at every time}\label{sec:unitary-entropy-every-time}

\paragraph{Files.}\leavevmode
\begin{itemize}
\item File: \lean{briefs/UnitaryEntropy/UnitaryEntropy/TwoParameter.lean}
\item Imports: \lean{SchroedingerEvolution.TwoParameter}, from the brief ``Multiple equivalent
  characterizations of the Schr\"odinger equation'' (for the two-parameter \lean{IsEvolution} and
  \lean{DrSceq} with $H(t)$; it imports this brief's main file), and the new
  \lean{ForPhyslib/QuantumMechanics/Wigner/Propagator.lean} (sections of determinant one, their
  uniqueness up to a root of unity, and gluing along an interval).
\item New \lean{ForMathlib} file \lean{ForMathlib/Analysis/Calculus/ContDiff/FiniteDimension.lean}:
  the determinant is a smooth function of the operator (Mathlib has only its continuity).
  \lean{ForMathlib/Analysis/Calculus/Deriv/Star.lean} gains
  \lean{star_mul_star_eq_neg_of_hasDerivAt_of_mem_unitary}: for a path of unitaries,
  $U'(t)U(t)^\dagger$ is skew-adjoint at every $t$. \lean{Wigner.exists_unitary_differentiableAt}
  now also gives differentiability at every $t$ near $t_0$ at which the family is differentiable.
\end{itemize}

\paragraph{Setting.} $E$ is finite-dimensional. A two-parameter evolution $\Phi_{t_0\to t_1}$,
$t_0 \le t_1$ (\lean{IsEvolution}), maps density operators affinely to density operators, with
$\Phi_{t\to t} = I$, $\Phi_{t_1\to t_2}\circ\Phi_{t_0\to t_1} = \Phi_{t_0\to t_2}$, and
$(\rho, t_0, t_1) \mapsto \Phi_{t_0\to t_1}(\rho)$ jointly continuous. There is no EVTI. DR-INFO
(\lean{DrInfo}) is $S(\Phi_{t_0\to t_1}(\rho)) = S(\rho)$ for all $t_0 \le t_1$. EVDIFF
(\lean{EvDiff}) is: for every reference time $r$ and density operator $\rho$,
$t \mapsto \Phi_{r\to t}(\rho)$ is differentiable at every $t > r$. DR-SCEQ (\lean{DrSceq}) is
conjugation, $\Phi_{t_0\to t_1}(\rho) = U(t_1)U(t_0)^\dagger\,\rho\,(U(t_1)U(t_0)^\dagger)^\dagger$, by
a $U$ with $U(0) = 1$ and $U'(t) = (i\hbar)^{-1}H(t)U(t)$ at every $t$, for self-adjoint $H(t)$.

\paragraph{Results.}
\begin{center}
\begin{tabularx}{\textwidth}{@{}BL@{}}
\toprule
Brief & Lean \\
\midrule
DR-INFO, EVDIFF for $\Phi_{t_0\to t_1}$ & \lean{DrInfo}, \lean{EvDiff} \\
Each $\Phi_{r\to t}$ is unitary & \lean{exists_unitary} \\
Local unitaries of determinant one & \lean{Wigner.eventually_isDetOneSection} \\
Uniqueness up to a constant root of unity & \lean{Wigner.IsDetOneSection.exists_eq_smul},
  \lean{Wigner.eq_inner_smul_of_proj_eq} \\
One differentiable $U(t)$ on an interval & \lean{Wigner.exists_isDetOneSection} (gluing:
  \lean{IsPreconnected.exists_forall_of_eventually}), \lean{exists_section} \\
One differentiable $U(t)$ for all $t$ & \lean{exists_propagator} \\
Step 3 at every time & \lean{star_mul_star_eq_neg_of_hasDerivAt_of_mem_unitary},
  \lean{exists_hamiltonian_of_star_eq_neg} \\
Theorem 2, at every time & \lean{drSceq_iff_drInfo_and_evDiff} \\
\bottomrule
\end{tabularx}
\end{center}

\paragraph{Differences from the brief.}
\begin{itemize}
\item Finite dimension only; see the next paragraph for why the argument does not extend.
\item EVDIFF asks for differentiability of $t \mapsto \Phi_{r\to t}(\rho)$ only at $t > r$,
  two-sided. There is no differentiability in the initial time and no one-sided derivative at
  $t = r$; a reference time $r < \tau$ gives the two-sided derivative of $U$ at $\tau$. A jump
  $H(t) = H_\beta$ for $t < 0$, $H(t) = H_\alpha$ for $t \ge 0$, with $H_\alpha - H_\beta$ not a
  multiple of $1$, violates EVDIFF at $t = 0$ (the one-sided derivatives
  $(i\hbar)^{-1}[H_{\alpha,\beta}, \rho(0)]$ differ for some $\rho$), as it must: the $U$ of
  DR-SCEQ is differentiable at every $t$.
\item DR-INFO has no reconstructed past; in finite dimension it is not needed.
\item The proof of ($\Leftarrow$). By Wigner's theorem each $\Phi_{r\to t}$ is induced by a unitary
  or an anti-unitary. Near any $t'$ the maps $\Phi_{\min(t,t')\to\max(t,t')}$ are close to the
  identity and so unitary (\lean{Wigner.eventually_exists_unitary}), and by the composition law
  $\Phi_{r\to t}$ has the type of $\Phi_{r\to t'}$; the type is locally constant, and
  $\Phi_{r\to r} = I$. Near $\tau > r$, conjugating by a unitary that induces $\Phi_{r\to\tau}$
  makes the family the identity at $\tau$, and the explicit \lean{Wigner.unitaryOf} is
  differentiable near $\tau$; multiplying by a branch of $\det(\cdot)^{-1/n}$, $n = \dim E$, gives
  determinant one. Two such families that induce the same maps differ by
  $c(t) = \braket{W(t)e|V(t)e}$, which is continuous with $c(t)^n = 1$, hence constant on a
  connected set. So the local families glue along $(r, \infty)$: the set of times reached from a
  fixed time is open and closed in $(r, \infty)$. The families for $r = -1, -2, \ldots$,
  normalized by $V(0)^\dagger$, agree, and give $U$; then $H(t) = i\hbar\,U'(t)U(t)^\dagger$.
\item The propagator $U_{t_0\to t_1} = U(t_1)U(t_0)^\dagger$ satisfies the composition law
  exactly; no phase cocycle appears.
\end{itemize}

\paragraph{Problems found in the brief.}
\begin{itemize}
\item The proof of Theorem~2 differentiates at one time, with phases chosen near that time.
  Going from ``at each time'' to one differentiable propagator for all times needs the local
  phase choices to be glued, and locally they are fixed only up to a continuous phase function.
  Fixing the determinant reduces this freedom to an $n$-th root of unity, which connectedness
  removes. The determinant needs finite dimension, so the infinite-dimensional case stays open.
\end{itemize}

\subsection{Hilbert spaces of any dimension}\label{sec:unitary-entropy-any-dim}

\paragraph{Files.}\leavevmode
\begin{itemize}
\item File: \lean{briefs/UnitaryEntropy/UnitaryEntropy/InfiniteDimensional.lean}
\item Library: \lean{UnitaryEntropy}, which now builds every module \lean{UnitaryEntropy.*}
\item Imports:\newline \lean{ForPhyslib.QuantumMechanics.VonNeumannEntropy},\newline
  \lean{ForPhyslib.QuantumMechanics.EqualMixtures} and \lean{ForPhyslib.QuantumMechanics.Kadison}
  (equal mixtures give affinity, and the mixed-state extension),\newline
  \lean{ForMathlib.Analysis.InnerProductSpace.StrongDerivative},\newline
  \lean{TauCeti.Analysis.InnerProductSpace.FunctionalCalculusComparison} (eigenvector action of
  the functional calculus, for the comparison in finite dimension),\newline
  and the finite-dimensional part \lean{UnitaryEntropy} (for \lean{born}, \lean{mixEntropy},
  \lean{mix_proj_cubic}, \lean{mix_proj_sq}, \lean{star_mul_self_of_schrodinger} and
  \lean{exists_hamiltonian}, none of which assumes finite dimension)
\item New file \lean{ForPhyslib/QuantumMechanics/DensityOperator.lean}: density operators on a
  Hilbert space of any dimension, on top of the trace-class operators of PhyslibAlpha
  (\lean{PhyslibAlpha/ProbabilisticTheory/HilbertSpace/TraceClass/}, from the physlib
  dependency; it builds against our Mathlib pin). Pure states, convex combinations, conjugation
  by an isometry, idempotent density operators are pure, and agreement with \lean{IsDensity} in
  finite dimension.
\item New file \lean{ForPhyslib/QuantumMechanics/VonNeumannEntropy.lean}: the von Neumann entropy
  in $[0,\infty]$, its invariance under conjugation by an isometry, and $S(\rho) = 0$ iff $\rho$
  is pure.
\item New \lean{ForMathlib} material:\newline
  \lean{ForMathlib/Analysis/InnerProductSpace/PositiveTrace.lean} (the trace of a positive
  operator in $[0,\infty]$, independent of the Hilbert basis),\newline
  \lean{ForMathlib/Analysis/InnerProductSpace/HilbertBasisMaps.lean} (for a Hilbert basis: a
  vector is determined by its coordinates, Parseval for their moduli, the diagonal unitary
  $b_i \mapsto \delta_i b_i$, coordinatewise conjugation, the isometry $b_i \mapsto v_i$ onto an
  orthonormal family, and $V(V^\dagger x) = x$ when $\|V^\dagger x\| = \|x\|$),\newline
  \lean{ForMathlib/Analysis/InnerProductSpace/StrongDerivative.lean}\newline (by Banach--Steinhaus,
  the strong derivative of an operator family is a bounded operator; for unitaries with
  $U(t_0) = 1$ it is skew-adjoint).
\item The Wigner library \lean{ForPhyslib/QuantumMechanics/Wigner/} (\lean{Basic},
  \lean{Normalize}, \lean{Rigidity}, \lean{Theorem}) now works over a Hilbert basis of any
  Hilbert space. \lean{Theorem} proves Wigner's theorem in any dimension, with isometries, and
  derives the finite-dimensional theorem from it. \lean{Continuity} excludes the anti-unitary
  branch in any dimension, without a trace, and derives the finite-dimensional exclusion from it.
  \lean{Differentiable} adds the differentiable choice of the phase in any dimension. The other
  files built on the library are unchanged.
\end{itemize}

\paragraph{Setting.} $E$ is a complex Hilbert space of any dimension; separability is not used.
\lean{posTrace}~$T = \sum_i \braket{e_i|Te_i} \in [0,\infty]$ for a Hilbert basis $(e_i)$. For
positive $T$ it does not depend on the basis (\lean{posTrace_eq_tsum}): write
$\braket{e_i|RR^\dagger e_i} = \|R^\dagger e_i\|^2$, expand by Parseval in a second basis and
exchange the two sums of nonnegative terms. A density operator (\lean{IsDensityOperator}) is a
positive operator that is trace class in the sense of PhyslibAlpha (\lean{IsTraceClass}) with
\lean{trace}~$\rho = 1$; equivalently, $0 \le \rho$ and \lean{posTrace}~$\rho = 1$
(\lean{isDensityOperator_iff}). The entropy is \lean{vonNeumannEntropy}~$\rho =$
\lean{posTrace}~$\eta(\rho)$ with $\eta(x) = -x\log x$ and $\eta(\rho)$ given by the continuous
functional calculus. The spectrum of a density operator lies in $[0,1]$, where $\eta \ge 0$, so
$\eta(\rho)$ is positive. \lean{proj}, \lean{mix}, \lean{born} and \lean{mixEntropy} are as above.

\paragraph{Results.}
\begin{center}
\begin{tabularx}{\textwidth}{@{}BL@{}}
\toprule
Brief & Lean \\
\midrule
QST:H, density operators & \lean{IsDensityOperator}, \lean{isDensityOperator_proj},
  \lean{IsDensityOperator.mix} \\
QST:H, von Neumann entropy & \lean{vonNeumannEntropy} \\
Agreement in finite dimension & \lean{isDensityOperator_iff_isDensity},
  \lean{vonNeumannEntropy_eq_ofReal_vonNeumann} \\
Corollary 1, entropy of $\tfrac12\rho_\psi + \tfrac12\rho_\phi$ &
  \lean{vonNeumannEntropy_mix_proj} \\
Corollary 1, $p(\phi|\psi) = \lambda$ iff $S = \ldots$ & \lean{born_eq_iff_vonNeumannEntropy}
  (for $\lambda \in [0,1]$) \\
Theorem 2 ($\Rightarrow$), unitary conjugation preserves the entropy &
  \lean{IsDensityOperator.vonNeumannEntropy_conj} (for isometries) \\
Theorem 2 ($\Rightarrow$), DR-INFO & \lean{IsDensityOperator.conj},
  \lean{schrodinger_conserves_vonNeumannEntropy} \\
Theorem 2 ($\Leftarrow$), Step 1 & \lean{IsDensityOperator.vonNeumannEntropy_eq_zero_iff} \\
Theorem 2 ($\Leftarrow$), Steps 1--2 up to Wigner &
  \lean{pure_and_born_of_preserves_vonNeumannEntropy}, \lean{exists_isWignerMap_of_preserves} \\
Wigner's theorem, any dimension &
  \lean{Wigner.IsWignerMap.exists_linearIsometry_or_conjLinearIsometry},
  \lean{Wigner.wigner_isometry} \\
Wigner's theorem, onto the rays &
  \lean{Wigner.IsWignerMap.exists_unitary_or_antiunitary_of_surjective} \\
Theorem 2 ($\Leftarrow$), Steps 1--2 with Wigner & \lean{exists_isometry_of_preserves};
  \lean{exists_unitary_of_preserves} if every pure state is an image \\
Theorem 2 ($\Leftarrow$), surjectivity from the reconstructed past & \lean{exists_proj_of_past},
  \lean{exists_unitary_of_preserves_of_past} \\
Theorem 2 ($\Leftarrow$), anti-unitary case excluded &
  \lean{eventually_exists_unitary_of_preserves} (from
  \lean{Wigner.eventually_exists_unitary_of_unitary_or_antiunitary}) \\
Theorem 2 ($\Leftarrow$), surjectivity is needed & \lean{shift_preserves_not_unitary} \\
Theorem 2 ($\Leftarrow$), differentiable $U(t)\psi$ & \lean{Wigner.exists_unitary_differentiableAt_apply} \\
Theorem 2 ($\Leftarrow$), Step 3 & \lean{ContinuousLinearMap.exists_hasDerivAt_apply},
  \lean{star_eq_neg_of_hasDerivAt_apply_of_mem_unitary}, \lean{exists_hamiltonian_of_star_eq_neg} \\
Theorem 2 ($\Leftarrow$), mixed states & \lean{affineOnDensityOperators_of_map_mix},
  \lean{AffineOnDensityOperators.eq_conjOp} \\
Theorem 2 ($\Leftarrow$), complete & \lean{schrodinger_of_preserves_vonNeumannEntropy} \\
\bottomrule
\end{tabularx}
\end{center}

\paragraph{Differences from the brief.}
\begin{itemize}
\item Wigner's theorem in any dimension follows Geh\'er's proof with a Hilbert basis $(b_k)$. The
  first normalization changes: the images $f(b_k)$ are orthonormal but need not span, so they
  define an isometry $V$ with $Vb_k = f(b_k)$, not a unitary. By Parseval,
  $\|V^\dagger f(v)\|^2 = \sum_k |\braket{f(b_k)|f(v)}|^2 = \sum_k |\braket{b_k|v}|^2 = 1 =
  \|f(v)\|^2$, so $f(v)$ lies in the range of $V$ and $V^\dagger \circ f$ is a Wigner map that
  fixes every $b_k$. The diagonal unitary and the conjugation are then built from the basis as in
  finite dimension. The result is a linear or conjugate-linear isometry; it is onto, hence
  unitary or anti-unitary, when every ray is the image of a ray, and always in finite dimension.
\item The anti-unitary branch is excluded by continuity without a trace. For $P = \ket{a}\bra{a}$,
  $Q$, $R$ rank-one projections onto $a, b, c$ and a fixed vector $x$,
  $\braket{x|PQRP\,x} = |\braket{a|x}|^2\,\Delta(a,b,c)$
  (\lean{inner_rankOne_mul_rankOne_mul_rankOne_mul_rankOne_apply}). Take the test triple
  $b_i, (b_i - b_k)/\sqrt2, (b_i - ib_k)/\sqrt2$ with $\Delta = (1+i)/4$, $x = b_i$, and the
  images $\Phi_t(P)$. The imaginary part is $\tfrac14$ at $t_0$ and continuous at $t_0$ in
  operator norm, but $\le 0$ for an anti-unitary, since its triple product is
  $\overline{\Delta}$ and the factor $|\braket{a|x}|^2$ is nonnegative. The finite-dimensional
  \lean{Wigner.eventually_exists_unitary}, which used $\operatorname{tr}(PQR) = \Delta$, is now
  a corollary. Continuity is assumed in operator norm; continuity in trace norm implies it.
\item EVDIFF is formalized as in finite dimension: $t \mapsto \Phi_t(\rho)$ is differentiable at
  $t_0$ in operator norm for every density operator $\rho$. The conclusion is the Schr\"odinger
  equation on vectors, as in DR-SCEQ: $\psi(t) = U(t)\psi$ satisfies $\psi'(t_0) = (i\hbar)^{-1}
  H\psi(t_0)$ for every $\psi$, with one bounded self-adjoint $H$. It is not the equation for
  $t \mapsto U(t)$ in operator norm, which the finite-dimensional \lean{schrodinger_of_preserves}
  gives; no choice of phases gives it (Section~\ref{sec:unitary-entropy-norm}). The basis sum that fixes the phase in finite dimension does not converge in operator
  norm; instead, for a fixed unit vector $e$, $\Phi_t$ extended by polarization gives
  $(\ket{U x}\bra{U e})\,e = \braket{U e|e}\,U x$, and dividing by $|\braket{U e|e}|$ fixes the
  phase one vector at a time. So $t \mapsto U(t)x$ is differentiable at $t_0$ for each $x$. By
  Banach--Steinhaus, applied to the difference quotients, the derivatives $x \mapsto
  (U(\cdot)x)'(t_0)$ form a bounded operator $A$. Differentiating $\braket{U(t)x|U(t)y} =
  \braket{x|y}$ gives $A^\dagger = -A$, and $H = i\hbar A$.
\item Mixed states follow as in finite dimension, with the spectral sum replaced by the canonical
  expansion $\rho = \sum_i \mu_i \ket{e_i}\bra{e_i}$, which may be infinite. Write
  $\rho = m\rho_F + (1-m)\sigma_F$ with $\rho_F$ a finite part of weight $m$; $\Phi_t$ and
  $U(t)\cdot U(t)^\dagger$ agree on $\rho_F$ by affinity, so they differ on $\rho$ by at most
  $2(1-m)$ in operator norm, and $m \to 1$ (\lean{AffineOnDensityOperators.eq_conjOp}).
  Affinity again comes from equal mixtures (\lean{affineOnDensityOperators_of_map_mix}).
\item No spectral decomposition is used. The entropy is the trace of $\eta(\rho)$, and pure states
  are found through the spectrum: $S(\rho) = 0$ forces $\eta(\rho) = 0$, so $\eta$ vanishes on
  the spectrum, which then lies in $\{0, 1\}$; $\rho$ is idempotent, and an idempotent density
  operator is $\ket{\psi}\bra{\psi}$ (\lean{IsDensityOperator.eq_proj_of_isIdempotentElem}).
\item Corollary~1 uses the cubic relation as in finite dimension, with the spectral mapping
  theorem for polynomials in place of eigenvalues. Then $\eta(\rho) = \alpha\rho + \beta\rho^2$,
  and its trace is computed from $\sum_i \braket{e_i|x}\braket{y|e_i} = \braket{y|x}$ for the
  rank-one terms.
\item The ($\Rightarrow$) direction uses only $U(t)^\dagger U(t) = 1$, which every propagator
  satisfies. Conjugation $X \mapsto UXU^\dagger$ by an isometry is a non-unital star algebra
  homomorphism, so it commutes with the functional calculus of $\eta$ (since $\eta(0) = 0$), and
  it preserves \lean{posTrace}. Unitarity is not needed.
\item Entropies are compared in $[0,\infty]$; in infinite dimension $S(\rho) = \infty$ occurs.
\item DR-INFO is formalized in full: $S(\Phi\rho) = S(\rho)$, and every density operator $\rho$
  has a reconstructed past, a density operator $\sigma$ with $\Phi\sigma = \rho$; then
  $S(\sigma) = S(\rho)$ by the first condition. The finite-dimensional part does not need the past,
  and \lean{exists_isometry_of_preserves} shows what remains without it.
\end{itemize}

\paragraph{Problems found in the brief.}
\begin{itemize}
\item In infinite dimension the reconstructed past in DR-INFO is essential, and Step~1 uses it
  when it says that pure states ``correspond bijectively''. DR-INFO presupposes a past ensemble
  $\rho_P$ for every initial ensemble, so $\Phi$ is onto the density operators, as in the
  introduction's ``deterministic and reversible''. Then every pure state is the image of a pure
  state, since its past has entropy zero (\lean{exists_proj_of_past}), and Wigner's theorem gives
  a unitary or an anti-unitary (\lean{exists_unitary_of_preserves_of_past}). The brief's proof
  is therefore correct, but the past cannot be dropped. Conjugation
  $\rho \mapsto S\rho S^\dagger$ by the shift $S$ on $\ell^2(\mathbb{N})$ maps density operators
  to density operators and preserves the entropy and equal mixtures, but no pure state goes to
  $\ket{e_0}\bra{e_0}$. So it is induced by no unitary and by no anti-unitary
  (\lean{shift_preserves_not_unitary}), and Wigner's theorem gives only an isometry. In finite
  dimension the past is automatic.
\item The time structure of [B]EVMIX does not replace the past. On $L^2(0,\infty)$ the right
  shifts $(V_tf)(x) = f(x-t)$ for $x > t$ and $0$ otherwise are isometries forming a strongly
  continuous semigroup with $V_0 = I$. Then $\Phi_{t_0 \to t_1}(\rho) =
  V_{t_1-t_0}\rho V_{t_1-t_0}^\dagger$ satisfies [B]EVMIX, preserves the entropy, and is induced
  by no unitary for $t_1 > t_0$; it has no reconstructed past. (By the Wold--Cooper
  decomposition, every strongly continuous semigroup of isometries is a unitary group plus copies of
  this one.) The example is formalized as \lean{shift_semigroup_counterexample} in the brief
  ``Multiple equivalent characterizations of the Schr\"odinger equation''; the Wold--Cooper
  decomposition is not. A norm-continuous semigroup of isometries is $V_t = e^{tA}$ with
  bounded $A$, and differentiating $V_t^\dagger V_t = I$ at $0$ gives $A^\dagger = -A$, so it is
  unitary. The example therefore needs an unbounded generator, which Step~3 does not cover
  anyway.
\item EVDIFF for every density operator forces a bounded Hamiltonian, as the Banach--Steinhaus
  step shows. Unbounded Hamiltonians, whose evolution is differentiable only on a dense set of
  states, are outside the theorem as stated; see Section~\ref{sec:unitary-entropy-unbounded}.
  Otherwise no problems beyond those of the finite-dimensional part.
\end{itemize}

\subsection{No norm-differentiable propagator}\label{sec:unitary-entropy-norm}

\paragraph{Files.}\leavevmode
\begin{itemize}
\item File: \lean{briefs/UnitaryEntropy/UnitaryEntropy/NormDifferentiable.lean}
\item Imports: \lean{ForPhyslib.QuantumMechanics.VonNeumannEntropy} and
  \lean{ForPhyslib.QuantumMechanics.Kadison} (density operators, entropy, products of rank-one
  projections) and \lean{ForPhyslib.QuantumMechanics.Wigner.Propagator} (two unitaries with the
  same action on pure states differ by a phase, \lean{Wigner.eq_inner_smul_of_proj_eq}).
\end{itemize}

\paragraph{Setting.} In any dimension, \lean{schrodinger_of_preserves_vonNeumannEntropy} chooses
the phases of the unitaries $U(t)$ so that $t \mapsto U(t)\psi$ is differentiable at $t_0$ for
every $\psi$ (Section~\ref{sec:unitary-entropy-any-dim}). In finite dimension the same
hypotheses make $t \mapsto U(t)$ differentiable in operator norm. In infinite dimension no choice
of phases does. On $\ell^2(\mathbb{N})$ with basis $(e_k)$ let
\[ U(t) = 1 + (e^{it} - 1)\,\ket{e_{k(t)}}\bra{e_{k(t)}}, \qquad
  k(t) = \lceil 1/|t| \rceil + 1, \]
which multiplies the component along $e_{k(t)}$ by $e^{it}$ (\lean{movingPhase}), and
$\Phi_t(\rho) = U(t)\rho\,U(t)^\dagger$. At $t = 0$ the factor $e^{it} - 1$ vanishes, so
$U(0) = 1$.

\paragraph{Results.}
\begin{center}
\begin{tabularx}{\textwidth}{@{}BL@{}}
\toprule
Statement & Lean \\
\midrule
$1 + (z-1)\ket{e}\bra{e}$ is unitary for $|z| = 1$ & \lean{phaseOp_mem_unitary} \\
$\|U\rho U^\dagger - \rho\| \le (2|z-1| + |z-1|^2)\,\|\rho e\|$ for self-adjoint $\rho$ &
  \lean{norm_phaseOp_mul_mul_star_sub_le} \\
$\|\rho e_n\| \to 0$ for a density operator $\rho$ & \lean{tendsto_norm_apply_hilbertBasis} \\
EVDIFF at $0$ in operator norm, with derivative $0$ & \lean{hasDerivAt_movingPhase_conj} \\
No unitary family inducing $\Phi_t$ near $0$ is norm-differentiable at $0$ &
  \lean{not_exists_differentiableAt_movingPhase} \\
All hypotheses of \lean{schrodinger_of_preserves_vonNeumannEntropy} hold, and no
  norm-differentiable $U(t)$ exists & \lean{movingPhase_counterexample} \\
\bottomrule
\end{tabularx}
\end{center}

\paragraph{Proof.}
\begin{itemize}
\item $\Phi_t$ is conjugation by a unitary, so it maps density operators to density operators,
  preserves the entropy and equal mixtures, and $\sigma = U(t)^\dagger\rho\,U(t)$ is a
  reconstructed past.
\item EVDIFF at $0$: with $P = \ket{e_k}\bra{e_k}$ and $a = e^{it} - 1$,
  $\Phi_t(\rho) - \rho = a\,P\rho + \bar a\,\rho P + |a|^2 P\rho P$, and $\|P\rho\| =
  \|\rho P\| = \|\rho e_k\|$. So $\|\Phi_t(\rho) - \rho\| \le (2|a| + |a|^2)\|\rho e_{k(t)}\|$,
  where $|a|/|t| \to 1$ and $\|\rho e_{k(t)}\| \to 0$ as $t \to 0$: by
  $\|\rho e_n\| \le \|\sqrt\rho\|\,\|\sqrt\rho\, e_n\|$ and
  $\sum_n \|\sqrt\rho\, e_n\|^2 = \operatorname{tr}\rho = 1$. The derivative is $0$ in operator
  norm, for every density operator.
\item No phase choice: if unitaries $V(t)$ induce $\Phi_t$ on pure states near $0$, then
  $V(t) = c(t)U(t)$ with $c(t) = \braket{e_0|V(t)e_0}$ and $|c(t)| = 1$, since
  $U(t)e_0 = e_0$. If $V$ is norm-differentiable at $0$, so are $c$ and $U = c^{-1}V$. For each
  $n$, $U(t)e_n = e_n$ near $t = 0$, so the derivative $U'(0)$ vanishes on every $e_n$, hence
  $U'(0) = 0$. But $\|U(t) - 1\| \ge \|(U(t) - 1)e_{k(t)}\| = |e^{it} - 1|$, and
  $|e^{it} - 1|/|t| \to 1$, so $U(t) - 1$ is not $o(t)$.
\end{itemize}

\paragraph{Notes.}
\begin{itemize}
\item The equation still holds on vectors: $U(t)\psi$ has derivative $0$ at $0$ for every
  $\psi$, so $H = 0$ in \lean{schrodinger_of_preserves_vonNeumannEntropy}. What fails is only the
  operator-norm form $U'(t_0) = (i\hbar)^{-1}HU(t_0)$, which the finite-dimensional
  \lean{schrodinger_of_preserves} gives.
\item EVDIFF holds here only at $t_0 = 0$: $k(t)$ jumps at $t = \pm 1/n$, and so does $U(t)$. A
  smooth variant (not formalized) puts the phase $\theta_k(t) = \varphi(kt)/k$, for a smooth bump
  $\varphi$ supported in $[1, 2]$ with $\varphi(3/2) \ne 0$, on every $e_k$ at once. Near each $t \ne 0$ only finitely
  many $\theta_k$ are nonzero, so $U(t)$ is norm-differentiable there; at $0$,
  $\|\Phi_t(\rho) - \rho\| \le 2\,\|U(t) - 1\|\,\|\rho Q_t\|$ with $Q_t$ the projection onto
  $\overline{\operatorname{span}}\{e_k : k \ge 1/|t|\}$, and $\|\rho Q_t\| \to 0$ since $\rho$ is
  compact. So EVDIFF holds at every time, also for the two-parameter family
  $U(t_1)U(t_0)^\dagger$, while $\|U(t) - 1\| \ge |e^{i\varphi(3/2)/k} - 1|$ at $t = 3/(2k)$, which
  is again not $o(t)$.
\item With EVTI the propagator is $e^{t(i\hbar)^{-1}H}$ with bounded $H$, which is
  norm-differentiable; the example needs a time-dependent evolution.
\end{itemize}

\paragraph{Problems found in the brief.}
\begin{itemize}
\item Step~3 differentiates $U(t)$ as an operator. In infinite dimension EVDIFF, even in
  operator norm on every density operator, does not give a norm-differentiable $U(t)$, only a
  strongly differentiable one; Step~3 then goes through on vectors, as in
  Section~\ref{sec:unitary-entropy-any-dim}.
\end{itemize}

\subsection{Unbounded Hamiltonians}\label{sec:unitary-entropy-unbounded}

\paragraph{Files.}\leavevmode
\begin{itemize}
\item File: \lean{briefs/UnitaryEntropy/UnitaryEntropy/Unbounded.lean}
\item Imports: \lean{SchroedingerEvolution.InfiniteDimensional}, the any-dimension part of the
  brief ``Multiple equivalent characterizations of the Schr\"odinger equation'' (for
  \lean{IsEvolution}, \lean{DrInfo}, \lean{DrSceq} and \lean{schrodinger_tfae}). That module
  imports \lean{UnitaryEntropy.InfiniteDimensional}, so the module graph stays acyclic, but the
  library \lean{UnitaryEntropy} now depends on the library \lean{SchroedingerEvolution}.
\end{itemize}

\paragraph{Setting.} As in Section~\ref{sec:unitary-entropy-any-dim}, restricted to a
time-independent evolution (EVTI): $\Phi_t$ for $t \ge 0$, with $\Phi_0 = I$,
$\Phi_{s+t} = \Phi_s \circ \Phi_t$, affine on the density operators and continuous at $0^+$ in
operator norm on each density operator (\lean{IsEvolution}). DR-INFO is entropy preservation
together with a reconstructed past for every density operator (\lean{DrInfo}). DR-SCEQ is
conjugation by a propagator $U$ with $U(0) = I$ that solves
$\tfrac{d}{dt}U(t)\psi = (i\hbar)^{-1}HU(t)\psi$ for $\psi$ in the domain of a self-adjoint,
in general unbounded, $H$ (\lean{DrSceq}). EVDIFF is weakened to a dense set of states
(\lean{EvDiff}): there is a dense subspace $D$ such that $t \mapsto \Phi_t(\ket\psi\bra\psi)$ is
differentiable on $t \ge 0$, in operator norm, for every unit vector $\psi \in D$.

\paragraph{Results.}
\begin{center}
\begin{tabularx}{\textwidth}{@{}BL@{}}
\toprule
Brief & Lean \\
\midrule
EVDIFF on a dense set of states & \lean{EvDiff} \\
Theorem 2 ($\Rightarrow$), EVDIFF & \lean{evDiff_of_drSceq} (with
  \lean{differentiableAt_proj_of_hasDerivAt}) \\
Theorem 2 with EVTI, unbounded $H$ & \lean{drSceq_iff_drInfo_and_evDiff} \\
\bottomrule
\end{tabularx}
\end{center}

\paragraph{Differences from the brief.}
\begin{itemize}
\item Only the time-independent case is formalized. For ($\Leftarrow$), DR-INFO gives DR-SCEQ
  by \lean{schrodinger_tfae}, through DR-PROB, DR-EV, DR-UNIT and Stone's theorem; EVDIFF is not
  used, since continuity and the group law already give a self-adjoint generator. For
  ($\Rightarrow$), DR-INFO is part of \lean{schrodinger_tfae}, and the dense subspace in EVDIFF
  is the domain of $H$: for $\psi$ there, $t \mapsto U(t)\psi$ is differentiable, hence so is
  $t \mapsto \ket{U(t)\psi}\bra{U(t)\psi} = \Phi_t(\ket\psi\bra\psi)$, as
  $(x, y) \mapsto \ket{x}\bra{y}$ is a bounded real-bilinear map.
\item The time-dependent case of Theorem~2 with unbounded $H(t)$ is not formalized; see the next
  paragraph.
\end{itemize}

\paragraph{Problems found in the brief.}
\begin{itemize}
\item For unbounded generators, Step~3 gives only a symmetric $H(t_0)$. If $t \mapsto U(t)\psi$
  is differentiable at $t_0$ only for $\psi$ in a dense subspace $D$, differentiating
  $\braket{U(t)\psi|U(t)\phi} = \braket{\psi|\phi}$ gives $\braket{A\psi|\phi} =
  -\braket{\psi|A\phi}$ on $D$: $A(t_0)$ is skew-symmetric, which is weaker than
  $A(t_0) = -A(t_0)^\dagger$ for an unbounded $A(t_0)$. Then $H(t_0)$ need not be self-adjoint,
  even on the largest such $D$. Example (not formalized): on $L^2(0,1)$ let $H_\alpha$ and
  $H_\beta$ be $-i\,d/dx$ with the boundary conditions $\psi(1) = e^{i\alpha}\psi(0)$ and
  $\psi(1) = e^{i\beta}\psi(0)$, $\alpha \ne \beta$, both self-adjoint. Take
  $H(t) = H_\beta$ for $t < 0$ and $H(t) = H_\alpha$ for $t \ge 0$, so the family
  differentiated at $t_0 = 0$ is $U(t) = U_{0\to t} = e^{-itH_\alpha/\hbar}$ for $t \ge 0$ and
  $U(t) = U_{t\to 0}^\dagger = e^{-itH_\beta/\hbar}$ for $t \le 0$.
  The evolution is unitary, so it satisfies [B]EVMIX and DR-INFO. The right derivative of
  $U(t)\psi$ at $0$ exists exactly for $\psi$ in the domain of $H_\alpha$, the left one for
  $\psi$ in the domain of $H_\beta$, so $U(t)\psi$ is differentiable at $0$ exactly when
  $\psi \in H^1(0,1)$ with $\psi(0) = \psi(1) = 0$. This subspace is dense, but $-i\,d/dx$ on it
  is the minimal momentum operator, closed and symmetric with deficiency indices $(1,1)$, and not
  self-adjoint. So the self-adjointness of $H(t)$ has to be assumed, as in Kato's theory of
  time-dependent Hamiltonians, rather than derived from EVDIFF. With EVTI it is derived, by
  Stone's theorem.
\end{itemize}

\end{document}
